\documentclass[journal]{IEEEtran}

\usepackage[utf8]{inputenc}
\usepackage[T1]{fontenc}
\usepackage{amsmath,amssymb}
\usepackage{graphicx}
\usepackage{booktabs}
\usepackage{multirow}
\usepackage{array}
\usepackage{url}
\usepackage{cite}
\usepackage{siunitx}
\usepackage{xcolor}
\usepackage{balance}
\usepackage{enumitem}
\usepackage{tabularx}

\sisetup{
  group-separator = {\,},
  detect-all
}

\newcommand{\Fp}{F_\mathrm{p}}
\newcommand{\Hz}{\,\text{Hz}}
\newcommand{\dB}{\,\text{dB}}

\begin{document}

\title{Measuring Koechlin's Timbral Attributes:\\Volume, Density, and Homogeneity\\in the Orchestral Spectrum}

\author{Yan~Maresz%
\thanks{Y.~Maresz is a composer and researcher at the CNSMDP, Paris. Correspondence: \texttt{y.maresz@cnsmdp.fr}. Data and code: \texttt{github.com/zseramnay/Formants}}%
}

\markboth{Référence Formantique Orchestrale --- Companion Paper~II}%
{Maresz: Measuring Koechlin's Timbral Attributes}

\maketitle

% ============================================================
\begin{abstract}
Charles Koechlin's \textit{Traité de l'orchestration} (1954) introduced a systematic theory of orchestral timbre centered on empirical attributes---volume, intensity, transparency, density---that predict how instruments blend into unified or heterogeneous orchestral planes. Despite the recognized originality of this framework, no attempt has been made to measure these attributes with modern spectral analysis tools. 

We propose a composite \textbf{Volume Index} combining four spectral descriptors---normalized spectral spread, sub-1\,kHz energy ratio, inverse formant dispersion, and MFCC spectral slope---computed across 5\,835 sustained samples from 53~orchestral instruments including muted variants. The index achieves a Spearman correlation of $\rho = -0.76$ ($p < 0.002$) with Koechlin's empirical volume scale. Crossing the resulting homogeneity matrix ($53 \times 53$) with formant convergence data ($\Delta F_1$, $\Delta\Fp$) per register yields a classification of 1\,378 instrument pairs into \textit{fondu} (blended), \textit{semi-fondu}, and \textit{hétérogène} (heterogeneous) planes, with 891~predicted fusions. Cross-register analysis reveals 1\,220 formant convergences distributed across five vowel zones from /u/ to /i/, confirming and extending the central convergence cluster previously identified in the 377--506\Hz{} band.
\end{abstract}

\begin{IEEEkeywords}
Orchestration, timbre, spectral analysis, Koechlin, formants, timbral fusion, orchestral planes, MFCC
\end{IEEEkeywords}

% ============================================================
\section{Introduction}
\label{sec:intro}

The question of why certain orchestral doublings produce a blended, homogeneous sound while others remain perceptually distinct has occupied orchestrators for over a century. Rimsky-Korsakov's \textit{Principes d'orchestration} (1914) offered instrument-specific descriptions and intensity equivalence rules, but no generalizable framework for predicting timbral fusion. As Chiasson and Duchesneau have shown~\cite{Chiasson2008, DuchesneauChiasson}, Charles Koechlin's \textit{Traité de l'orchestration}~\cite{Koechlin1954} represents a decisive advance: Koechlin proposed a coherent system of \emph{general attributes}---volume, intensity, transparency, density---applicable across all instruments, enabling systematic comparison of timbres and prediction of orchestral blend.

Koechlin's \textit{volume} is not loudness (intensity), but the apparent \emph{space occupied by a sound}: a low flute note is voluminous but weak, while a medium oboe note is thin but intense. This distinction, which Koechlin illustrated with an astronomical metaphor---the flute as a giant star of minimal density, the oboe as a dense dwarf star---anticipates S.~S.~Stevens' psychoacoustic work on tonal volume~\cite{Stevens1965} and more recent research on auditory size perception~\cite{VanDinther2006}. Yet as Chiasson notes~\cite{ChiassonHumaniste}, no experimental verification of Koechlin's volume scale across orchestral families has been attempted, and ``the universality of volume as an attribute of instrumental timbre remains to be confirmed.''

This paper proposes such a measurement. Building on the \textit{Référence Formantique Orchestrale}---a corpus of 5\,914+ spectral analyses from the Orchidea/SOL2020 database---we construct a composite spectral index that operationalizes Koechlin's concept of volume, then use it to build a homogeneity matrix predicting which instrument combinations will form blended orchestral planes.

% ============================================================
\section{Koechlin's Timbral Framework}
\label{sec:koechlin}

Following the analysis of Chiasson~\cite{Chiasson2008, ChiassonHumaniste} and Duchesneau \& Chiasson~\cite{DuchesneauChiasson}, Koechlin's system can be summarized as a hierarchy of attributes operating at three levels:

\textbf{Instrument level.} Each instrument possesses, at each register, a characteristic \textit{volume} (apparent size) and \textit{intensity} (perceived force). These vary independently: the flute's volume is large in the low register but its intensity is weak; the oboe increases in both intensity and volume toward the low register. From these two, a third attribute emerges: \textit{transparency} (large volume, low intensity = transparent; small volume, high intensity = dense).

\textbf{Orchestral plane level.} When instruments share equivalent volume \emph{and} intensity, they fuse into a single orchestral plane. This plane may exhibit \textit{fondu} (blend) if the timbres are sufficiently similar, and \textit{homogénéité} if the blend is particularly strong (as in a string section). Unequal volumes or intensities produce \textit{hétérogénéité}: multiple distinct planes heard simultaneously.

\textbf{Orchestration level.} The full orchestration achieves \textit{plénitude} (fullness)---either \textit{matérielle} (from voluminous, blended timbres with doublings) or \textit{d'écriture} (from dense counterpoint). \textit{Éclat} (brilliance) results from pure, undoubled timbres.

Koechlin provides an empirical scale of mean instrumental volumes (vol.~I, p.~288), ranging from ``gros'' (large: tubas, horns, trombones) to ``mince'' (thin: oboe, clarinet in the upper register). Our objective is to measure this scale spectrally and extend it to the full modern orchestra, including muted instruments.

% ============================================================
\section{Corpus and Descriptors}
\label{sec:corpus}

\subsection{Instrumental Corpus}

The analysis covers \textbf{53~instruments} (or instrument--mute combinations) totaling \textbf{5\,835 sustained samples} across \textbf{193~register segments}. Sources are the Orchidea/SOL2020 database (35~entries) and the Yan\_Adds complementary database (18~entries). The corpus includes:

\begin{itemize}[nosep]
\item \textbf{Flutes:} concert, bass, contrabass, piccolo
\item \textbf{Double reeds:} oboe (open, muted), English horn
\item \textbf{Single reeds:} clarinet B$\flat$, E$\flat$, bass, contrabass; saxophone alto
\item \textbf{Bassoons:} bassoon (open, muted), contrabassoon
\item \textbf{Brass:} horn (open, muted), trumpet~C (open, cup, harmon, straight, wah-open, wah-closed), trombone (same 5~variants), bass trombone, bass tuba (open, muted), contrabass tuba
\item \textbf{Strings solo:} violin, viola, cello, bass---each open, with standard mute, and with lead mute (\textit{sordina piombo})
\item \textbf{String ensembles:} violin, viola, cello, bass sections, plus muted variants
\item \textbf{Plucked strings:} guitar, harp
\item \textbf{Pitched percussion:} marimba, vibraphone
\end{itemize}

Register segmentation follows the definitions in \texttt{registres.md}, with muted variants inheriting the registers of their base instrument. For brass mutes, only \textit{ordinario} samples are retained (for wah mutes: \textit{ordinario\_open} and \textit{ordinario\_closed}).

\subsection{Spectral Descriptors}

Four descriptors are combined into the Volume Index, each capturing a different facet of the perceived ``size'' of a sound:

\textbf{$V_1$: Normalized Spectral Spread} $= \sigma_\text{spectral} / \mu_\text{centroid}$. The spectral spread (second moment, standard deviation of the spectral distribution) divided by the spectral centroid (first moment). This ratio captures the \emph{relative} spectral width, normalized for pitch. Instruments with broad harmonic energy relative to their centroid---tubas, low strings---have high $V_1$.

\textbf{$V_2$: Sub-1\,kHz Energy Ratio} $= E_{<1\text{kHz}} / E_\text{total}$. Computed from the 1024-bin linear amplitude spectrum (FFT = 4096, $f_s = 44100$\Hz). The proportion of squared spectral energy below 1\,kHz captures Koechlin's intuition that voluminous instruments concentrate energy in the low frequencies. The threshold of 1\,kHz was chosen as the approximate boundary between the fundamental/low-harmonic region and the formant/brilliance region.

\textbf{$V_3$: Inverse Formant Dispersion} $= 1 / \overline{\Delta F}$, where $\overline{\Delta F}$ is the mean spacing between consecutive formants $F_1$ through $F_6$ (from the formant extraction pipeline described in the companion paper~\cite{Maresz2025}). This descriptor is grounded in auditory size perception research~\cite{Fitch1997, VanDinther2006}: closely-spaced formants signal a large resonating body, widely-spaced formants a small one. $V_3$ is constant per instrument (formants are extracted globally, not per register).

\textbf{$V_4$: MFCC$_1$ (Spectral Slope).} The first Mel-Frequency Cepstral Coefficient (after $C_0$), computed from a 26-band mel filterbank applied to the amplitude spectrum. MFCC$_1$ captures the overall \emph{tilt} of the spectral envelope: positive values indicate energy concentrated in the low mel bands (``warm'' timbres), negative values indicate energy shifted toward high frequencies (``bright'' timbres). This complements $V_2$ by operating on a perceptually-scaled frequency axis.

\subsection{Composite Volume Index}

Each descriptor $V_k$ is $z$-score normalized across all 193~register entries:
\begin{equation}
\hat{V}_k = \frac{V_k - \mu_k}{\sigma_k}
\end{equation}

The composite Volume Index is the arithmetic mean of the available normalized components:
\begin{equation}
\mathcal{V} = \frac{1}{|\mathcal{K}|} \sum_{k \in \mathcal{K}} \hat{V}_k
\end{equation}
where $\mathcal{K} \subseteq \{1,2,3,4\}$ is the set of descriptors available for the given instrument--register combination.

\textbf{Density Index.} Following Stevens' formulation that loudness is the product of volume and density~\cite{Stevens1965}, at constant intensity: $\mathcal{D} = -\mathcal{V}$.

% ============================================================
\section{Results: Volume Index}
\label{sec:volume}

\subsection{Ranking and Comparison with Koechlin}

Table~\ref{tab:volume_top} shows the Volume Index for the reference register of each instrument family.

\begin{table}[!t]
\centering
\caption{Volume Index $\mathcal{V}$ at reference register (médium or equivalent). Koechlin rank: 1=gros, 6=mince.}
\label{tab:volume_top}
\small
\begin{tabular}{@{}lccrc@{}}
\toprule
\textbf{Instrument} & \textbf{Reg.} & $\mathcal{V}$ & \textbf{K} \\
\midrule
Contrabass Ensemble     & grave   & $+1.49$ & 1 \\
Contrabass (sordina)    & grave   & $+1.36$ & 2 \\
Contrabass              & grave   & $+1.32$ & 2 \\
Contrabass Tuba         & médium  & $+1.22$ & 1 \\
Bass Tuba               & médium  & $+1.19$ & 1 \\
Bass Flute              & grave   & $+1.22$ & -- \\
Trombone (harmon)       & grave   & $+1.19$ & -- \\
Contrabass Flute        & grave   & $+1.18$ & -- \\
Horn                    & médium  & $+0.55$ & 2 \\
Trombone                & médium  & $+0.56$ & 3 \\
Bassoon                 & grave   & $+0.60$ & 4 \\
Violoncello             & médium  & $+0.48$ & 4 \\
Flute                   & médium  & $-0.06$ & 4 \\
Clarinet B$\flat$       & chalu.  & $+0.23$ & 5 \\
Viola                   & médium  & $-0.21$ & 6 \\
Violin                  & médium  & $-0.73$ & 6 \\
Oboe                    & médium  & $-0.92$ & 5 \\
Trumpet C               & médium  & $-0.64$ & 4 \\
Piccolo                 & grave   & $-1.10$ & -- \\
Trumpet C (cup)         & grave   & $-1.36$ & -- \\
Trumpet C (harmon)      & grave   & $-1.33$ & -- \\
\bottomrule
\end{tabular}
\end{table}

The Spearman rank correlation between $\mathcal{V}$ and Koechlin's empirical scale is $\rho = -0.76$ ($p < 0.002$, $n = 17$ base instruments). The negative sign is expected: high $\mathcal{V}$ corresponds to Koechlin's rank~1 (``gros''). The ordering reproduces Koechlin's scale almost exactly, with two noteworthy exceptions.

\subsection{The Trumpet Anomaly}

The trumpet~C, which Koechlin ranks alongside bassoons and flutes (rank~4, ``medium''), falls to $\mathcal{V} = -0.64$ in our analysis. All four components contribute to this low score: its spectral energy is concentrated above 1\,kHz ($V_2 = 0.47$), its spread/centroid ratio is low ($V_1 = 1.05$), and its MFCC$_1$ is modest ($V_4 = 32.0$). This suggests that Koechlin's perception of the trumpet's ``volume'' may have incorporated a spatial projection quality---a capacity to fill a concert hall---that our spectral descriptors, computed from close-microphone recordings, do not capture. Room acoustics and directivity may be the missing factor.

\subsection{Effects of Mutes}

Mutes produce dramatic shifts in the Volume Index. Table~\ref{tab:mutes} summarizes the effect.

\begin{table}[!t]
\centering
\caption{Effect of mutes on the Volume Index.}
\label{tab:mutes}
\small
\begin{tabular}{@{}lcc@{}}
\toprule
\textbf{Combination} & $\Delta\mathcal{V}$ & \textbf{Effect} \\
\midrule
Trumpet + cup        & $-0.72$ & thins dramatically \\
Trumpet + harmon     & $-0.69$ & thins dramatically \\
Trumpet + straight   & $-0.35$ & moderate thinning \\
Trombone + harmon    & $+0.63$ & \textit{increases} volume \\
Horn + sordina       & $-0.08$ & negligible \\
Strings + sordina    & $\approx 0$ & slight darkening \\
Strings + piombo     & $-0.3$ to $-0.8$ & strong thinning \\
\bottomrule
\end{tabular}
\end{table}

The most striking result is the trombone harmon mute, which \emph{increases} the Volume Index by $+0.63$. The harmon mute suppresses upper harmonics while reinforcing low-frequency energy, producing a spectrally ``fatter'' sound. This contrasts sharply with the trumpet harmon, which does the opposite---a finding consistent with the asymmetric mute effects documented in the companion paper.

% ============================================================
\section{Homogeneity Matrix}
\label{sec:homogeneity}

\subsection{Construction}

The homogeneity between two instruments $A$ and $B$ is computed from their reference-register data as:
\begin{equation}
H(A, B) = 1 - \alpha \cdot \frac{|\mathcal{V}_A - \mathcal{V}_B|}{\mathcal{V}_\text{range}} - \beta \cdot d_\text{cos}(\mathbf{m}_A, \mathbf{m}_B)
\end{equation}
where $\mathcal{V}_\text{range}$ is the total Volume Index range across all instruments, $d_\text{cos}$ is the cosine distance between the MFCC profiles (coefficients 1--12, omitting $C_0$), and $\alpha = \beta = 0.5$.

The MFCC cosine distance captures similarity of spectral \emph{shape} independently of overall energy level, complementing the Volume Index which captures spectral \emph{size}. The resulting $53 \times 53$ matrix encodes Koechlin's concept of homogeneity: $H \geq 0.85$ indicates instruments likely to form a blended (``fondu'') plan; $H < 0.70$ indicates heterogeneous timbres.

\subsection{Key Findings}

The highest cross-family homogeneities in the matrix confirm established orchestral practice while revealing unexpected pairings:

\textbf{Horn + Trombone} ($H = 0.99$): the most homogeneous cross-family pair, confirming the archetypal Brahms/Wagner brass blend. Their MFCC profiles are nearly superimposed, and their Volume Indices differ by only 0.01.

\textbf{Bass Clarinet + Horn} ($H = 0.99$), \textbf{Bass Clarinet + Trombone} ($H = 0.99$): the bass clarinet emerges as a spectral ``bridge'' between woodwinds and brass, a finding that explains its extensive use in the French orchestral tradition (Debussy, Ravel).

\textbf{Horn + Marimba} ($H = 0.96$): an unexpected result---these instruments share nearly identical F1 values at the medium register and similar spectral profiles. This pairing, rarely exploited in classical orchestration, is extensively used in contemporary music.

\textbf{Trumpet + Violin} ($H = 0.98$): at the medium register, these instruments share similar density profiles and MFCC shapes, confirming the classical doubling of the melodic line.

\textbf{Oboe: the outsider.} The oboe shows consistently low homogeneity with brass and low strings ($H < 0.70$), confirming Koechlin's description of it as a ``thin, solid'' timbre---a ``dwarf star of considerable density.'' Its spectral profile is distinct from all other instruments, which is precisely why it stands out so effectively as a solo voice.

% ============================================================
\section{Orchestral Planes: Crossing Volume with Formant Convergence}
\label{sec:planes}

\subsection{Methodology}

The homogeneity matrix captures similarity of spectral \emph{profile}, but Koechlin's theory requires an additional condition for blend: the instruments must also be in the same acoustic ``space.'' We operationalize this through formant convergence: two instruments form a blended plane when their principal formants ($F_1$ and $\Fp$) are close \emph{and} their spectral profiles are similar.

For each of the $\binom{53}{2} = 1\,378$ instrument pairs, we compute:
\begin{itemize}[nosep]
\item $\Delta F_1$: absolute difference of the dominant spectral peak (80--1500\Hz) at the reference register, extracted from the mean spectral envelope
\item $\Delta \Fp$: absolute difference of the energy-weighted spectral centroid in the 800--1800\Hz{} band
\item A \textit{fusion score}: $0.4 \cdot H + 0.3 \cdot (1 - \Delta F_1/500) + 0.3 \cdot (1 - \Delta\Fp/500)$
\end{itemize}

Pairs are classified as:
\begin{itemize}[nosep]
\item \textbf{FONDU} ($\star$): $H \geq 0.80$ and ($\Delta F_1 \leq 30$ or $\Delta\Fp \leq 50$\Hz)
\item \textbf{SEMI-FONDU} ($\bullet$): $H \geq 0.70$ and $\Delta F_1 \leq 80$, or $H \geq 0.80$ alone
\item \textbf{H\'ET\'EROG\`ENE} ($\circ$): remaining pairs
\end{itemize}

\subsection{Distribution of Predicted Planes}

\begin{table}[!t]
\centering
\caption{Classification of 1\,378 instrument pairs.}
\label{tab:planes}
\begin{tabular}{@{}lrc@{}}
\toprule
\textbf{Category} & \textbf{N pairs} & \textbf{\%} \\
\midrule
FONDU (F1 convergence)      & 148 & 10.7\% \\
FONDU (Fp convergence)      & 160 & 11.6\% \\
SEMI-FONDU (F1)             &  95 &  6.9\% \\
SEMI-FONDU (H only)         & 488 & 35.4\% \\
CONVERGENT (F1 only)        &  11 &  0.8\% \\
HÉTÉROGÈNE                  & 476 & 34.5\% \\
\bottomrule
\end{tabular}
\end{table}

The ``CONVERGENT'' category is particularly interesting: it captures 11~pairs where $\Delta F_1 \leq 30$\Hz{} but $H < 0.70$---instruments in the same ``vowel zone'' but with very different spectral volumes. In Koechlin's terms, these would form \textit{superimposed heterogeneous planes in the same acoustic color}: audible separately but sharing the same harmonic character.

\subsection{Notable Cross-Family Fusions}

Among the highest-scoring cross-family blended pairs (Table~\ref{tab:fusions}), several extend well beyond classical orchestration practice:

\begin{table}[!t]
\centering
\caption{Top cross-family predicted fusions (different instrument families).}
\label{tab:fusions}
\small
\begin{tabular}{@{}llccc@{}}
\toprule
\textbf{Instrument A} & \textbf{Instrument B} & $H$ & $\Delta F_1$ & $\Delta\Fp$ \\
\midrule
Cl.~Mib         & Flute            & 0.95 &  0  &  5 \\
Tuba (sord.)    & Basson (sord.)   & 1.00 & 11  & 30 \\
Horn            & Marimba          & 0.96 &  0  & 19 \\
Fl.~cb          & Alto (sord.)     & 0.98 & 11  & 23 \\
Fl.~basse       & Alto (sord.)     & 0.97 & 10  & 19 \\
Tuba            & Guitare          & 0.94 & 11  &  1 \\
Fl.~basse       & Vibraphone       & 0.92 &  0  &  -- \\
Cor (sord.)     & Vibraphone       & 0.95 &  0  &  -- \\
Tbn (harmon)    & Vc (piombo)      & 0.97 & 22  & 12 \\
Alto (sord.)    & Vn~ens.~(sord.)  & 0.95 & 22  &  -- \\
\bottomrule
\end{tabular}
\end{table}

The trombone harmon + cello piombo pairing ($\Delta\Fp = 12$\Hz) is a discovery: the harmon mute shifts the trombone spectrum into a zone that converges precisely with the lead-muted cello. Similarly, the bass tuba + guitar pairing ($\Delta\Fp = 1$\Hz) suggests that the guitar functions as a kind of ``plucked tuba'' in spectral terms---a relationship hinted at in the classical guitar-and-cello repertoire but never quantified.

% ============================================================
\section{Vowel Zone Analysis of Convergences}
\label{sec:vowels}

Extending the analysis to all register combinations, we identify 1\,220 instrument pairs with $\Delta F_1 \leq 30$\Hz. These convergences cluster into five vowel zones:

\begin{table}[!t]
\centering
\caption{Distribution of formant convergences by vowel zone.}
\label{tab:zones}
\begin{tabular}{@{}lcrl@{}}
\toprule
\textbf{Zone} & \textbf{F1 band} & \textbf{N} & \textbf{Orchestral character} \\
\midrule
/u/       & $\leq 250$\Hz    & 544 & Foundation; \textit{plénitude matérielle} \\
/o/       & 251--400\Hz      & 423 & Central convergence cluster \\
/\aa/     & 401--520\Hz      & 125 & Transitional warmth \\
/a/       & 521--800\Hz      & 54  & Presence, projection \\
/e/--/i/  & $> 800$\Hz       & 74  & Brilliance; muted brass \\
\bottomrule
\end{tabular}
\end{table}

\textbf{79\% of all convergences} occur in the /u/ and /o/ zones (F1 $\leq 400$\Hz), confirming that the foundation of orchestral blend lies in the low-frequency formant region. This is consistent with Koechlin's emphasis on the \textit{sous-médium} register (the octave below middle~C) as the locus of \textit{plénitude matérielle}.

The /o/ zone (251--400\Hz, 423~convergences) corresponds exactly to the convergence cluster identified in the companion paper at 377--506\Hz. It is here that the cross-family pairings central to the orchestral tradition converge: horn + viola ($\Delta F_1 = 11$\Hz), violin + bassoon ($\Delta F_1 = 11$\Hz), English horn + clarinet ($\Delta F_1 = 11$\Hz).

The /e/--/i/ zone ($> 800$\Hz, 74~convergences) is almost entirely populated by \emph{muted instruments}: trumpet cup, trumpet harmon, and piccolo. These instruments, whose mutes shift $F_1$ dramatically upward, create an ``upper convergence cluster'' that has no equivalent in classical orchestration theory.

% ============================================================
\section{Discussion}
\label{sec:discussion}

\subsection{Validating Koechlin Empirically}

The strong correlation between our spectral Volume Index and Koechlin's empirical scale ($\rho = -0.76$) suggests that Koechlin's auditory perception of instrumental ``volume''---developed through decades of orchestral practice without access to spectral analysis tools---was capturing real, measurable properties of the acoustic signal. The four descriptors we combine---normalized spectral width, low-frequency energy proportion, formant spacing, and spectral slope---jointly approximate a percept that Koechlin could only describe metaphorically.

\subsection{Beyond Koechlin: The Mute Revolution}

Koechlin's framework, developed for the open orchestral instruments of his era, gains unexpected power when extended to muted instruments. Mutes do not merely attenuate---they \emph{restructure} the spectral profile, creating new convergences that cross traditional family boundaries. The trombone harmon, which \textit{increases} spectral volume (unlike all other brass mutes), and the trumpet cup, which shifts $F_1$ into the /e/--/i/ zone, represent timbral transformations that no classical treatise has systematized.

\subsection{Limitations}

The Volume Index uses equal weighting ($\frac{1}{4}$ per descriptor), which may not reflect the perceptual salience of each component. Perceptual listening tests would be needed to calibrate the weights. Furthermore, the index is computed from close-microphone anechoic recordings, which do not capture the spatial radiation patterns that undoubtedly influence perceived volume in a concert hall---this likely explains the trumpet anomaly.

% ============================================================
\section{Conclusion}
\label{sec:conclusion}

We have demonstrated that Koechlin's empirical timbral attributes---volume, density, transparency, and homogeneity---can be operationalized through spectral analysis and extended to the full modern orchestra including muted instruments. The composite Volume Index, the $53 \times 53$ homogeneity matrix, and the cross-register convergence map constitute a quantitative framework for predicting orchestral blend that both validates and extends the most comprehensive theory of orchestration in the literature.

The data (7~CSV files and analysis scripts) are available at \texttt{github.com/zseramnay/Orchestration}.

% ============================================================
\section*{Acknowledgments}
This work is part of the \textit{Référence Formantique Orchestrale} project at the CNSMDP, Paris. Spectral data from the Orchidea/SOL2020 database (IRCAM) and the Yan\_Adds complementary corpus.

\begin{thebibliography}{99}

\bibitem{Koechlin1954}
C.~Koechlin, \textit{Traité de l'orchestration}, 4~vols., Paris: Max Eschig, 1954--1959.

\bibitem{Chiasson2008}
F.~Chiasson, ``L'originalité du Traité de Koechlin: étude comparative avec d'autres traités d'orchestration,'' Conference OICCM, Université de Montréal, 2008.

\bibitem{DuchesneauChiasson}
M.~Duchesneau and F.~Chiasson, ``Le Traité de l'orchestration de Charles Koechlin,'' in N.~Donin and L.~Feneyrou (eds.), \textit{Théorie et composition musicales au XX\textsuperscript{e} siècle}, Lyon: Symétrie, 2013.

\bibitem{ChiassonHumaniste}
F.~Chiasson, ``L'universalité de la méthode de Koechlin,'' in \textit{Koechlin, compositeur et humaniste}, Paris: Vrin, 2010.

\bibitem{Stevens1965}
S.~S.~Stevens, M.~Guirao, and A.~W.~Slawson, ``Loudness, a product of volume times density,'' \textit{J.~Experimental Psychology}, vol.~69, pp.~503--510, 1965.

\bibitem{VanDinther2006}
R.~van Dinther and R.~D.~Patterson, ``Perception of acoustic scale and size in musical instrument sounds,'' \textit{J.~Acoust.~Soc.~Am.}, vol.~120, no.~4, pp.~2158--2176, 2006.

\bibitem{Fitch1997}
W.~T.~Fitch, ``Vocal tract length and formant frequency dispersion correlate with body size in rhesus macaques,'' \textit{J.~Acoust.~Soc.~Am.}, vol.~102, no.~2, pp.~1213--1222, 1997.

\bibitem{Maresz2025}
Y.~Maresz, ``Formant analysis of orchestral instruments: a spectral reference for timbral convergence and classical doubling practices,'' \textit{Référence Formantique Orchestrale}, CNSMDP, 2025. Available: \texttt{github.com/zseramnay/Formants}.

\bibitem{Sandell1991}
G.~J.~Sandell, ``Concurrent timbres in orchestration: a perceptual study of factors determining blend,'' Ph.D.\ thesis, Northwestern University, 1991.

\bibitem{RimskyKorsakov1914}
N.~Rimsky-Korsakov and M.~Steinberg, \textit{Principes d'orchestration}, Berlin: Edition russe de musique, 1914.

\bibitem{Meyer2009}
J.~Meyer, \textit{Acoustics and the Performance of Music}, 5th ed., New York: Springer, 2009.

\end{thebibliography}

\end{document}
